% Section 10: the AXIOM LIST.
% RULE FOR THIS SECTION: an axiom entry appears here if and only if some proof
% in this document uses it and does not prove it. Every axiom entry is a closed
% statement with its hypotheses written out and is referenced by number from
% each consumer. The internal theorem stands separately from
% the imported interface.
\section{The axiom list}\label{sec:axioms}

This section contains the complete eight-entry interface of the document: the
statements it uses and does not prove. A formalization of the main theorem must
supply each of them, either by importing a library result or by proving it.
The eight axiom entries fall into two groups --- the three inputs declared by the
campaign, and five theorems taken from the literature. The document also proves
one internal cyclic-classification theorem, Theorem~\ref{thm:mycyclic}, proved in Section~\ref{sec:transport}, which stands separately from that interface.

\subsection{Declared inputs of the campaign}

\begin{axiom}[Hypothesis R for $X_1$]\label{ax:R}
$X_1(P_+)=X_1(P_-)$ across every simple Reidemeister III event, i.e.\ every
event whose zero set is the forced bundle
$Z=\{\mathrm{G3}_{e,f,g},\mathrm{G4}_{e;f,g},\mathrm{G4}_{f;e,g},
\mathrm{G4}_{g;e,f}\}$ for three pairwise remote edges with
$1\leq e<f<g\leq n$, concurrent at $t=0$ at a point interior to all three, the
event being transversal in the sense of Definition~\ref{def:event}.

\emph{Why it is not proved here.} It is not a theorem of the literature and not
a claim of this document: it is a \emph{declared input} of the campaign, the
hypothesis under which the main theorem is stated. Proving it is the open
problem the campaign exists for, and a document that proved it would not need
the entry.
\emph{Used in:} Theorem~\ref{thm:main}(B), and there only --- in the Reidemeister
III case of its induction. No other statement of this document reads it.
\end{axiom}

\begin{remark}[the domain of Hypothesis R, and what the entry said before]
\label{rem:axRhistory}
The index order and the transversality in Axiom~\ref{ax:R} are the antecedent
of (S2) in Definition~\ref{def:laws}, and that entry now carries the same one: a
declared input and the law it declares must have one domain, and an earlier
revision repaired the law's antecedent and left the entry's alone.

\emph{Fidelity: this domain is the specification's own.} The specification asks
$C(P_+)=C(P_-)$ across every Reidemeister III wall, and restricts every one of
its laws to a \emph{simple} event: only the predicate defining the named event
vanishes, all other listed vertex, edge-concurrency and crossing-order
predicates keeping their signs. Write $\mathcal E_{\mathrm{spec}}$ for the
events its (S2) thereby covers and $\mathcal E_{\mathrm{ax}}$ for those this
entry covers. Then $\mathcal E_{\mathrm{ax}}=\mathcal E_{\mathrm{spec}}$.
Containment one way is immediate, an event admitted here being a
Reidemeister-III wall at which nothing outside the forced bundle vanishes. The
other way is the reading of "the predicate defining the event": at a transverse
concurrency the two crossings on each of the three edges collide, so the three
order predicates vanish with it necessarily --- they are dependents of the
concurrency, not independent coincidences --- and a specification that asked
them to keep their signs there would have an empty (S2). The defining condition
is the bundle, exactly as the flat wall's three vanishing members are one
condition; and with that reading the two classes are the same events.

An earlier revision of this paragraph called the entry's domain a strict
subset and defended the narrowing as the safe direction for an assumption. That
was wrong in the direction that flatters: it claimed less than the entry
delivers, and it misread the specification, whose simplicity clause was already
doing the restricting. The correction is the peer's.

Where the law is \emph{read} is narrower than where it is asserted, and that is
a fact about the proof rather than about the entry: its one consumer,
Theorem~\ref{thm:main}(B), meets Reidemeister-III walls only at certified simple-factor hits of the
path supplied by Lemma~\ref{lem:relgp}(B), and it proves there that the zero set is
the forced bundle rather than assuming it here. A sentence in this position previously said that "nothing is lost
downstream by the restriction" and spoke of "the excluded events"; with the
equality above there are no excluded events, and the sentence was left over from
the framing the paragraph before it withdraws.
\emph{Wording changed, hypothesis unchanged.} An earlier revision of this entry
read "a concurrency of three pairwise remote edges together with the
crossing-order predicates it exchanges, all other members of $\mathcal G$
keeping their signs". That domain is \emph{empty}: at a transverse concurrency
the two crossings on each of the three edges collide, so all three order
predicates vanish with the concurrency and cannot be required to keep their
signs. \emph{And the change is not neutral.} A universal law whose domain is empty is
vacuously true, so the earlier print asserted nothing and the present entry,
whose domain is the actual Reidemeister-III zero set, is formally
\textbf{stronger}. What is unchanged is the \emph{intended} input: the
specification's own words, ``$C(P_+)=C(P_-)$ across every Reidemeister III
wall''. That intention was formalized vacuously in the earlier revision and is
formalized on the real domain here. An earlier revision of this paragraph
claimed the entry was ``neither strengthened nor weakened''; that was wrong.
\end{remark}



\begin{remark}[why the parameter is a convention and not a phrase]
\label{rem:lawfulbinder}
An earlier revision wrote the first entry as a closed existential, "there is a
function $A$ which \ldots", and then used the letter $A$ in the second entry and
in the theorem with nothing binding it. Two closed existentials have two
witnesses: the peer's two-object model satisfies each separately and neither
jointly, and the document contains no common binder for them.

A second revision --- also mine --- read $A$ as \emph{universally} quantified
and described that in print as strictly stronger than the existential form. It
is not: the two are \emph{incomparable}. With $X_1=0$ and no function satisfying
both properties, the universal statement is vacuously true and the existential
false; with two such functions, one agreeing with $X_1$ and one not, the
existential-and-agreement is true and the universal false. What the campaign
assumes is neither: its owner \emph{supplies} the quantity, so $A$ is a declared
constant (Convention~\ref{conv:lawful}) of which the two entries below assert
properties. The content assumed is unchanged by either repair; what moved is
where the quantifier sits, and it now has a numbered home a formalizer can point
at.

\emph{What the axiom entries of this section are.} Each axiom entry is an
assumption about objects already declared: about $X_1$, which
Definition~\ref{def:X1} constructs; about the declared constant $A$; or about
the mathematical literature, in which case it is a theorem quoted and not
proved here. No axiom entry declares an object of its own, and none asserts a
bare existence.
Theorem~\ref{thm:mycyclic} stands separately as a result proved in Section~\ref{sec:transport}.
\end{remark}

\begin{axiom}[the lawful quantity]\label{ax:lawful}
With $A$ as in Convention~\ref{conv:lawful}, $A$ is constant on chambers,
unchanged across silent events, and satisfies (S2), (S3), (S4), (S5), (S7) with
the sign (S6), (S8), and (S9) of Definition~\ref{def:laws}, each on the domain
where Definition~\ref{def:laws} states it. Write $\mathrm{Lawful}(A)$ for this
conjunction. In particular
\begin{equation}
  A(\sigma^{k}P)=A(P)\qquad\text{for every generic }P\text{ and every }k\geq0,
  \label{eq:lawful-cyclic}
\end{equation}
by induction on $k$ from \eqref{eq:S9}.

\emph{Why it is not proved here.} It is a declared input, of the same kind as
the rest of the entry: the properties the campaign's owner asserts of the
quantity $A$ it supplies. Nothing in this document computes $A$.

\emph{Stability of the remaining laws under \eqref{eq:sigma}.} The relabeling
permutes the labels and fixes the point set, the plane and its orientation. The
census is stated as equalities evaluated at $\sigma P$, in the direction of
\eqref{eq:pullback}.

\begin{itemize}
\item[(chambers)] The relabeling is a homeomorphism of the generic locus with
  inverse $\sigma^{n-1}$, by Definition~\ref{def:cycshift}, so it carries
  chambers to chambers and the constancy clause transports.
\item[(silence)] Definition~\ref{def:silent} asks of \emph{every} member of the
  zero set $Z$ that it be a predicate $\mathrm{G2}_{e,i}$ whose vanishing at the
  event places $p_i$ on the \emph{line} of $e$ but not on the segment $e$; in
  particular no member of $Z$ is a turn (G1), a concurrency (G3), a
  crossing-order predicate (G4) or a direction determinant (G5). Both halves
  transport. By \eqref{eq:pullback} the $\mathrm{G2}$ expressions satisfy
  $\mathrm{G2}_{e,i}(\sigma P)=\mathrm{G2}_{e^{+},i^{+}}(P)$, and the
  correspondence of index tuples is a bijection of each guard family with
  itself, so $Z(\sigma P)$ consists of $\mathrm{G2}$ members exactly when $Z(P)$
  does and the exclusion of the other four families is preserved. For the
  remaining half, the line of an edge and the segment spanned by it are
  determined by the edge as a point set; $\sigma$ fixes every edge as a point
  set and moves no vertex, so the corresponding member's vanishing places the
  same vertex on the same line and off the same segment. The condition holds at
  $\sigma P$ for every member of $Z(\sigma P)$ exactly when it holds at $P$ for
  every member of $Z(P)$, and silent events correspond to silent events.
\item[(genericity)] The guards are indexed by labels, and evaluation at
  $\sigma P$ reads the guard at the successor tuple:
  $\mathrm{G1}_i(\sigma P)=\mathrm{G1}_{i^{+}}(P)$ and
  $\mathrm{G2}_{e,i}(\sigma P)=\mathrm{G2}_{e^{+},i^{+}}(P)$, and likewise for
  $\mathrm{G3}$, $\mathrm{G4}$ and $\mathrm{G5}$ on their index tuples.
  Remoteness, incidence and the crossing of two segments are conditions on the
  edges and vertices rather than on their labels, and $\sigma$ fixes those point
  sets, so the activation conditions of $\mathrm{G3}$, $\mathrm{G4}$ and
  $\mathrm{G5}$ hold at $\sigma P$ exactly at the tuples corresponding under
  \eqref{eq:pullback}. The generic locus is carried to itself.
% UNSIGNED TYPESETTING (D5-ACC-3): item-heading repair; rival re-verification pending.
\item[(cut)]\emph{Guard values at the cut.} The conditional members carry an
  \emph{ordered representative}: (G5) is indexed by $e<f$, (G3) by $e<f<g$ and
  (G4) by a distinguished edge together with $f<g$, the comparisons being
  between the integer representatives in $\{1,\dots,n\}$, as
  Definition~\ref{def:guarded} prescribes. The relabeling sends the edge with
  representative $i$ at $\sigma P$ to the edge with representative $i^{+}$ at
  $P$, and $i\mapsto i^{+}$ preserves the order of every pair except one: the
  top representative $n$ wraps to $1$. The three identities, each an equality of
  values \emph{evaluated at} $\sigma P$, are therefore
  \begin{align}
    \mathrm{G5}_{e,f}(\sigma P)&=
      \begin{cases}
        \mathrm{G5}_{e^{+},f^{+}}(P), & f<n,\\[2pt]
        -\,\mathrm{G5}_{1,e^{+}}(P), & f=n,
      \end{cases}
    \label{eq:g5-wrap}\\[4pt]
    \mathrm{G4}_{e;f,g}(\sigma P)&=
      \begin{cases}
        \mathrm{G4}_{e^{+};f^{+},g^{+}}(P), & g<n,\\[2pt]
        -\,\mathrm{G4}_{e^{+};1,f^{+}}(P), & g=n,
      \end{cases}
    \label{eq:g4-wrap}\\[4pt]
    \mathrm{G3}_{e,f,g}(\sigma P)&=
      \begin{cases}
        \mathrm{G3}_{e^{+},f^{+},g^{+}}(P), & g<n,\\[2pt]
        \phantom{-\,}\mathrm{G3}_{1,e^{+},f^{+}}(P), & g=n.
      \end{cases}
    \label{eq:g3-wrap}
  \end{align}
  Off the wrap the representatives increase with the labels and the member is
  read unchanged. At the wrap the two determinantal families behave differently,
  and the difference is the parity of the permutation that restores the sorted
  order. In \eqref{eq:g5-wrap} with $f=n$ the pair $(e^{+},f^{+})=(e^{+},1)$ is
  \emph{decreasing}, and $\det(d_f,d_e)=-\det(d_e,d_f)$ gives the sign. In
  \eqref{eq:g4-wrap} with $g=n$ the ordered pair $(f^{+},1)$ is likewise
  decreasing, and the antisymmetry
  $\mathrm{G4}_{e;g,f}=-\mathrm{G4}_{e;f,g}$ of Definition~\ref{def:guarded}
  gives the sign; the distinguished edge carries no order constraint and simply
  becomes $e^{+}$. In \eqref{eq:g3-wrap} with $g=n$ the three representatives
  present themselves as $(e^{+},f^{+},1)$ and the sorted triple is
  $(1,e^{+},f^{+})$, reached by a cyclic permutation of three rows; that
  permutation is \emph{even}, so the $3\times3$ determinant is unchanged and no
  sign appears. Genericity asks that the guard values be nonvanishing, which
  every case above preserves; the two sign changes are changes of the
  representative, not of the geometry.
\item[(S2)] The domain asks that the zero set be the forced bundle
  $\{\mathrm{G3}_{e,f,g},\mathrm{G4}_{e;f,g},\mathrm{G4}_{f;e,g},
  \mathrm{G4}_{g;e,f}\}$ of three pairwise remote edges, concurrent and
  transversal. All four members correspond under \eqref{eq:pullback} to the
  members of the successor triple, remoteness and concurrency are conditions on
  the edges, and transversality is the simultaneous sign change of the bundle,
  which the permutation preserves. The asserted equality $C(P_+)=C(P_-)$ carries
  no label.
\item[(S3), (S4)] The events and the genericity of the wall deletion are
  conditions on the polygon. The statements evaluate $C$ on the deletion, and
  \begin{equation}
    (\sigma P)\setminus j^{-}=
    \begin{cases}
      P\setminus j, & j=1,\\[2pt]
      \sigma\,(P\setminus j), & j\neq1,
    \end{cases}
    \label{eq:deletion-shift}
  \end{equation}
  as configurations. In the case $j=1$ the deleted label crosses the cut and the
  child is the deletion itself; otherwise the child is a relabeling of it, and
  $C$ agrees on the two by \eqref{eq:S9} applied at arity $n-1$. The left and
  right labels of the wall, which belong to (S3), and the loop and no-loop
  labels together with the change $\kappa$ of the Whitney rotation number, which
  belong to (S4), are read from the event's geometry and are carried with it,
  $\sigma$ fixing that geometry. Equation \eqref{eq:deletion-shift} is an
  identity of configurations, and it is at that level that it is used.
\item[(S5)] The empty-spike condition is a condition on the disc bounded by the
  loop-side arc, and the asserted value $0$ carries no label.
\item[(S6), (S7)] The sign $s$ of \eqref{eq:S6} is the sign of
  $\det(p_b-p_a,\,p_M-p_a)$, formed from the oriented remote edge $(a,b)$ and
  the moving vertex $M$ and read in the orientation of the plane; $\sigma$ moves
  the labels $a,b,M$ and fixes the three points and the orientation, so $s$ is
  unchanged. The halves of Definition~\ref{def:halves} are based at $M$ and
  correspond under \eqref{eq:pullback} to the halves of the relabeled polygon,
  so
  \begin{equation}
    C(P_+)-C(P_-)=s\,C(\lambda_1)C(\lambda_2)
    \label{eq:lawful-s7}
  \end{equation}
  transports with the same sign, the values $C(\lambda_i)$ agreeing with those
  of the relabeled halves by \eqref{eq:S9} at the halves' arities.
\item[(S8)] Strict convexity is stated in the cyclic order of the hull boundary,
  and equivalently by the turns having one sign with no three vertices collinear
  and the curve embedded; each condition is invariant under \eqref{eq:sigma}.
  The two asserted values, $1$ for the clockwise boundary and $(-1)^{n}$ for the
  counterclockwise one, are attached to the two orientations of the boundary,
  which $\sigma$ preserves, so each convex polygon keeps the value of its own
  orientation.
\end{itemize}

\emph{Used in:} the consumers of the clauses already present are unchanged ---
Theorem~\ref{thm:zeroanchor}(B) reads (S5), and Theorem~\ref{thm:main}(B) reads
all of them, additionally reading \eqref{eq:S9} at the relabeled endpoint of its
transport step. Those two consume it. Four
remarks --- Remarks~\ref{rem:noncircular}, \ref{rem:notassumed},
\ref{rem:s7wellposed} and \ref{rem:silencenew} --- name it while discussing what
is assumed and what is proved; naming is not consuming.
\end{axiom}

\begin{remark}[the transport read of \eqref{eq:S9}]
\label{rem:lawful-cyclic-use}
The transport step of Theorem~\ref{thm:main}(B) terminates at
$\sigma^{k}Z_{n,r}$, where $Z_{n,r}$ is the anchor selected there and $k$ is the
relabeling supplied by Theorem~\ref{thm:mycyclic}. The step reads
\begin{equation}
  A\bigl(\sigma^{k}Z_{n,r}\bigr)=A\bigl(Z_{n,r}\bigr)
  \label{eq:anchor-shift}
\end{equation}
for that $k$ and that anchor. Equation \eqref{eq:anchor-shift} is read for every
iterate returned by Theorem~\ref{thm:mycyclic} and at every selected anchor,
which is the generality \eqref{eq:lawful-cyclic} supplies.

The declared anchor values depend on the rank alone and \eqref{eq:sigma}
preserves the rank, so anchors related by \eqref{eq:S9} carry equal declared
values. At $(n,r)=(4,0)$ the anchor is $K_0$ and \eqref{eq:anchor-shift} carries
the value $A(K_0)=-1$ of Axiom~\ref{ax:ownerlist} to each $\sigma^{k}K_0$.
\end{remark}

\begin{axiom}[owner normalization on the anchors]\label{ax:ownerlist}
With $A$ as in Convention~\ref{conv:lawful}, $\Cat$ the Catalan numbers of
Definition~\ref{def:catalan} --- the display below reads $\Cat_{|r|}$, the
subscript bound there by $r$ --- and $K_r$ the circular star of
Definition~\ref{def:circularstar} for an integer $r$ with $|r|\geq1$, together
with $K_0$ the four-vertex bow-tie of that definition:
\[
   A(K_0)=-1,\qquad
   A(K_r)=\sgn(r)\,(-1)^{|r|}\,\Cat_{|r|}\quad\text{for every integer }r\neq0 .
\]
The quantifier is written out --- $r$ ranging over the nonzero integers ---
because an earlier revision left it to be read off the display; the Catalan
numbers' own quantifier lives with their definition.
Write $\mathrm{Owner}(A)$ for this pair of identities.

\emph{Why it is not proved here.} A declared input again: the two normalization
values the owner supplies on the anchors. This document computes $X_1$ on those
anchors --- that is Proposition~\ref{prop:anchorvalues}(D) --- but it does not and
cannot compute $A$ there.
\emph{Used in:} Theorem~\ref{thm:main}(B), which reads the anchor values
\emph{directly} --- at the base case $n=3$, where $\Delta(K_{\pm1})=0$ needs
$A(K_{\pm1})$, and at the minimal anchors of the induction step. An earlier
revision routed the use "through the rank-zero anchor of
Theorem~\ref{thm:zeroanchor}(B)", whose proof reads Axiom~\ref{ax:lawful}'s (S5)
clause and not this entry. That is the one consumer. Remarks
\ref{rem:notassumed}, \ref{rem:s7wellposed} and \ref{rem:whichanchors} name it
in discussion, which is not the same as reading it.
\end{axiom}

\begin{remark}[an entry that was added and then struck]\label{rem:diffdeadremoved}
An intermediate revision of this document listed a fourth entry in this group,
asserting that at a noninterlacing bigon wall an eligible support whose two
newborns lie in different residual pieces has a mixed contact carrier. It was
added because the different-piece row had been reduced to a single term whose
prefactor is that selector, with nothing yet fixing that selector's value.

It was false. The peer produced an exact guarded $24$-gon with a live
selector on a different-piece support, recorded with its data in
Remark~\ref{rem:diffcensus}. The entry is struck, and the sector it was
covering is now proved without it in Theorem~\ref{thm:s7universal}(D)(v): the row dies
by the coefficient, not by the selector. The episode is left in the text
because an axiom introduced to cover a gap in a proof is exactly the move
this document is supposed to make visible.
\end{remark}

\subsection{Literature axioms}

\begin{axiom}[HOMFLY--PT]\label{ax:homfly}
There is a unique map $L\mapsto P_L(a,z)\in\mathbb Z[a^{\pm1},z^{\pm1}]$ from
isotopy classes of oriented links in $S^3$ satisfying
$P_{\bigcirc}=1$ and $aP_{L_+}-a^{-1}P_{L_-}=zP_{L_0}$ for every skein triple.
In particular $P_L$ is invariant under the Reidemeister moves. One consequence
is consumed as part of this entry: $P_K\in\mathbb Z[a^{\pm1},z^2]$ for a
\emph{knot} $K$, so that the $z^0$ coefficient of a product of knot
polynomials is the product of their $z^0$ coefficients.

\emph{Why it is not proved here.} The existence of the invariant is a theorem
about all links at once, proved by induction over the full skein tree with a
descending complexity. This document does not construct that invariant; the relative induction in
Lemma~\ref{lem:homflyrows} starts with $P$ already supplied and compares
based-link expressions only.
\emph{Used in:} Definition~\ref{def:X1} (to have $P_H$ at all);
Proposition~\ref{prop:chamberinv}(ii), Lemma~\ref{lem:flatdata} and
Lemma~\ref{lem:silence}; Lemma~\ref{lem:curl} and Theorem~\ref{thm:carrierfloor},
clauses~(C) and~(R); Lemma~\ref{lem:homflyrows},
Proposition~\ref{prop:sectortarget}(B), Lemma~\ref{lem:fulltwist},
Lemma~\ref{lem:triplebridge} and Lemma~\ref{lem:nitransport}, clause~(A); and, through the
knot-parity clause, Theorem~\ref{thm:s7universal}(D)(i) and
Theorem~\ref{thm:s7universal}(D)(iv)(c).
\end{axiom}

\begin{remark}[the trimmed clauses of Axiom~\ref{ax:homfly}]
\label{rem:homflywhole}
This remark is about one entry and two clauses, and it makes no claim about
any other.

\emph{Axiom~\ref{ax:homfly}, clause dropped.} The disjoint-union value
$P_{L\sqcup\bigcirc}=\frac{a-a^{-1}}{z}P_L$ left the entry as a duplicate and
not as a narrowing: it is a consequence of the skein relation and it is
displayed in Definition~\ref{def:homfly} where the normalization is fixed.

\emph{Axiom~\ref{ax:homfly}, clause still dropped.} The general
$z$-order bound for a multi-component link is not restored.
Lemma~\ref{lem:homflyrows}(ii) uses only the retained knot clause: smoothing a
mixed crossing gives a knot, and
$[z^{-2}]P_K=0$ because $P_K\in\mathbb Z[a^{\pm1},z^2]$ has nonnegative even
$z$-powers. It then proves the two-component coefficient identity directly by
mixed switches and split separation. No order bound for a multi-component
polynomial is assumed.

\end{remark}

\begin{axiom}[Etnyre; self-linking from a front]\label{ax:etnyre}
Let $T$ be a front diagram of a transverse knot with no downward vertical
tangency. Then $\operatorname{sl}(T)$ equals the writhe of $T$.

\emph{Why it is not proved here.} The identity is a computation of the
self-linking number from a front, which presupposes the contact-geometric
definition of $\operatorname{sl}$; this document never defines it, and reads
the identity only as a bridge from the diagram's writhe to the transverse
HOMFLY--PT bound of Axiom~\ref{ax:slbound}.
\emph{Used in:} Theorem~\ref{thm:carrierfloor}(C).
\end{axiom}

\begin{axiom}[the HOMFLY--PT bound on the self-linking number]\label{ax:slbound}
Let $T$ be a transverse knot in the standard contact $\mathbb R^3$ and let
$P_T(a,z)$ be the HOMFLY--PT polynomial of the knot it presents, in the
normalization of Axiom~\ref{ax:homfly}. Then
\[
   \operatorname{sl}(T)\;\leq\;-\max\deg_a P_T(a,z)-1 .
\]

\emph{Why it is not proved here.} The bound is a theorem about the whole
invariant, proved in the source by a skein template argument this document has
no machinery for. Source fidelity: the source states the bound for the maximal
self-linking number of a topological knot type and proves it representative by
representative, each front's self-linking number being that of the usual
positive transverse pushoff; the form assumed here, at a single transverse
representative $T$, is that maximum read at one of its values, and that
reading is quoted rather than renamed: every transverse knot is the positive
transverse pushoff of a Legendrian knot, by the Legendrian-pushoff
construction of Etnyre's survey, Section~2.9, and the positive pushoff's
self-linking number is $\operatorname{tb}-r$ of that Legendrian by Lemma~2.22
of the same section.  So $\operatorname{sl}(T)$ is one of the values the
maximum is taken over. The normalization of the source is that of Axiom~\ref{ax:homfly}
verbatim, in the same variable $a$; no variable conversion is assumed.
\emph{Used in:} Theorem~\ref{thm:carrierfloor}(C).
\end{axiom}

\begin{axiom}[Brini--Eynard--Mari\~no; one coefficient of a torus knot]
\label{ax:bem}
For every integer $r\geq2$, with $T(r,2r+1)$ the torus knot of those parameters,
$P$ the HOMFLY--PT polynomial in the normalization of
Definition~\ref{def:homfly}, and $\Cat_r$ the Catalan number of Definition~\ref{def:catalan},
\[
   \bigl[a^{\,1-(2r+1)(r-1)-r}\,z^{0}\bigr]\,P_{T(r,2r+1)}(a,z)
   =(-1)^{r-1}\,\Cat_r .
\]

\emph{Used in:} Proposition~\ref{prop:anchorvalues}(D), and there only.
\emph{Why it is not proved here.} The coefficient is read off a closed
evaluation of the torus-knot polynomials obtained from topological recursion on
a spectral curve; nothing in this document computes a torus-knot polynomial.
Remark~\ref{rem:bemfidelity} carries that evaluation, the variable dictionary
and the derivation of the identity above from it, so the step from the
literature to this entry is on the page.
\end{axiom}

\begin{remark}[the torus-knot entry and its source formula]
\label{rem:bemfidelity}
\emph{The source's statement, in full.} Let $P,Q$ be coprime with
$P\geq Q\geq1$ --- coprime so that $K_{Q,P}$ is a knot, and $P\geq Q$ so that
every factorial below has a nonnegative argument at every $k\leq d$ --- and put
$d=\min(P-1,Q-1)=Q-1$, reading $[0]!=1$. With $q$ the quantum variable,
\[
   [n]=q^{n/2}-q^{-n/2},\qquad [n]!=[n][n-1]\cdots[1],\qquad
   z=[1]=q^{1/2}-q^{-1/2},
\]
and $c=e^{t/2}$ the exponentiated 't~Hooft coupling, the HOMFLY--PT polynomial
of the torus knot $K_{Q,P}$ in the normalization of
Definition~\ref{def:homfly}, written in the campaign variables $a=c^{-1}$ and
$z=q^{1/2}-q^{-1/2}$, satisfies
\begin{equation}
   \frac{P_{K_{Q,P}}(a,z)}{z}
   =c^{(P-1)(Q-1)}\,\frac{1}{[P][Q]}
   \sum_{k=0}^{d}(-1)^k c^{2k}\,
   \frac{[P+Q-k-1]!}{[P-k-1]!\,[Q-k-1]!\,[k]!} .
\label{eq:bemsum}
\end{equation}
That display is the source's own, read at \texttt{arxiv.org/abs/1105.2012}:
quantum integers as displayed, $c=e^{t/2}$, the sum to $d=\min(P-1,Q-1)$, and
the left side normalized by $q^{1/2}-q^{-1/2}$. Axiom~\ref{ax:bem} assumes one
coefficient of it and not the display, and the derivation of that coefficient
from this one is below, so the step from this display to what the entry
assumes is on the page, next to its source.

\emph{Why the entry is the coefficient and not the sum.} Only one number of
\eqref{eq:bemsum} is ever read, and the step that used the rest --- that the
terminal summand has the least $a$-degree in the row --- is needed only to
extract that number \emph{from} the sum. Stating the number directly removes
the step; an earlier revision kept the sum as the entry and defended it by that
step.

For the row the consumer extracts, the bookkeeping of $z$-powers has to be
exact, and an earlier revision of this paragraph was off by one. Since
$[n]=nz+O(z^3)$, a quantum factorial $[m]!$ has $z$-order $m$, so the summand
$[P+Q-k-1]!/\bigl([P-k-1]!\,[Q-k-1]!\,[k]!\bigr)$ has $z$-order
$(P+Q-k-1)-(P-k-1)-(Q-k-1)-k=1$, while $[P][Q]$ has $z$-order $2$: the right
side, and so the left, has $z$-order $-1$. The row the consumer wants is
therefore the coefficient of $z^{-1}$ in $P_{K_{Q,P}}/z$, which is the
coefficient of $z^0$ in $P_{K_{Q,P}}$ itself:
\[
   \bigl[z^{-1}\bigr]\bigl(P_{K_{Q,P}}/z\bigr)
   =\bigl[z^{0}\bigr]P_{K_{Q,P}}
   =\sum_k(-1)^k c^{(P-1)(Q-1)+2k}\,C_k/(PQ),
   \qquad
   C_k=\frac{(P+Q-k-1)!}{(P-k-1)!\,(Q-k-1)!\,k!} .
\]
Writing $[z^0]$ of the divided left side, as that revision did, names a
coefficient one power too high. For $(Q,P)=(2,3)$ this gives
$2c^{2}-c^{4}=2a^{-2}-a^{-4}$, the $z^0$ row of the positive trefoil; for
$(Q,P)=(2,5)$ it gives $3c^{4}-2c^{6}=3a^{-4}-2a^{-6}$, the $z^0$ row of the
$(2,5)$ torus knot. \emph{The derivation of Axiom~\ref{ax:bem} from \eqref{eq:bemsum}, in five
steps.} (1) The $z^0$ row of the polynomial is the displayed sum's constant
term, by the bookkeeping just given. (2) The $k$-th summand carries the campaign
exponent $-\bigl((P-1)(Q-1)+2k\bigr)$ in $a$, since $c=a^{-1}$, and that
exponent is strictly decreasing in $k$; so the summands sit at distinct
$a$-degrees and the terminal one, $k=Q-1$, sits at the least. (3) Its quantum
ratio has classical limit
$C_{Q-1}=\frac{P!}{(P-Q)!\,(Q-1)!}$, so its coefficient is
$(-1)^{Q-1}\frac1Q\binom{P-1}{Q-1}$ at $a^{1-P(Q-1)-Q}$. (4) Specialize to the
pairs the circular stars present, $(Q,P)=(r,2r+1)$ with $r\geq2$: the exponent
is $1-(2r+1)(r-1)-r$ and the coefficient is
$(-1)^{r-1}\frac1r\binom{2r}{r-1}$. (5) That number is $(-1)^{r-1}\Cat_r$,
since $\frac1r\binom{2r}{r-1}$ and $\frac1{r+1}\binom{2r}{r}$ are both
$\frac{(2r)!}{r!\,(r+1)!}$. Steps (1)--(5) are exactly the identity
Axiom~\ref{ax:bem} states, and step (2) --- the least-degree step --- lives
here, where it is used to read one number off a formula, and nowhere in the
document's own chain.
\end{remark}

\begin{remark}[an entry removed]\label{rem:prremoved}
The phase-1.5 manuscript proved the genericity of the circular star by
importing Poonen and Rubinstein's theorem on the diagonals of a regular
$n$-gon. Lemma~\ref{lem:stargeneric}(i) proves it instead from the fact that the
star's edges are tangent to one concentric circle and that at most two tangents
to a circle pass through a point, so that import is gone from this list. Its
removal is the kind of change the axiom list exists to make visible.
\end{remark}

\begin{axiom}[isomorphic records present the same link]\label{ax:gausscode}
Let $D$ and $D'$ be oriented link diagrams, each of whose underlying plane
curves is a connected generic immersed circle in the sphere --- one component,
transverse double points, no triple points --- and suppose there is an
orientation-preserving record isomorphism from the record of $D$ to the record
of $D'$ (Definition~\ref{def:record}). Then $D$ and $D'$ present the same
oriented link.
\emph{Used in:} Lemma~\ref{lem:pieceintrinsic}, Lemma~\ref{lem:curl},
Proposition~\ref{prop:chamberinv}(ii), Lemma~\ref{lem:flatdata},
Lemma~\ref{lem:silence}, Corollary~\ref{cor:groupedknot}(A) and~(B),
Proposition~\ref{prop:sectortarget}(B),
Lemma~\ref{lem:slidingcartesian},
clauses~(A) and~(C), Theorem~\ref{thm:s7universal}, clause~(A) and
Lemma~\ref{lem:triplebridge}. An earlier revision named the first and then said
``every consumer of a grouped carrier row'', which is a description and not a
list.
\emph{Why it is not proved here.} This is the classical statement that a
realizable classical signed Gauss code determines its oriented link. Its content
is Carter's classification of generic immersed curves by their Gauss data: two
such curves in a surface with isomorphic data differ by a homeomorphism of the
surface, \emph{provided the curves fill it}, that is, provided every
complementary region is a disc. The hypothesis above is the case in which that
proviso is automatic and is stated for that reason rather than dropped: the
underlying curve of each diagram is a connected four-valent graph embedded in
the sphere, and every face of a connected graph embedded in the sphere is a
disc, the crossing-free case being a circle with two disc faces. The resulting
homeomorphism carries diagram to diagram and so link to link, a priori only up
to reflection if it reverses orientation; the crossing signs, which the record
isomorphism preserves, exclude that. This document does not develop the
realizability theory, and the consumers need the conclusion for diagrams built
on different curves, whose traversal circles are compared by an explicit record
isomorphism rather than identified.
\emph{Two earlier revisions.} The first did not list this entry at all: the
identification it provides was asserted inside the proof of
Corollary~\ref{cor:groupedknot}(B) in one sentence, and the peer refuted the
sentence --- smoothing the support changes the traversal successor, so retaining
the same labels does not identify the cyclic signed records. The second listed
the entry but hypothesized that the two records were \emph{equal}, which
presupposes a shared traversal circle that these two diagrams do not have; the
peer refused that typing too. The combinatorial content is proved ---
Lemma~\ref{lem:carrierword} for the order, Lemma~\ref{lem:pieceintrinsic} for
the map --- and what remains, the passage from a record isomorphism to an
equality of links, is this entry.
\end{axiom}

\begin{remark}[an entry removed, the second]\label{rem:arcremoved}
An earlier revision carried an axiom asserting that two embedded regular arcs
in a disc with the same endpoints and endpoint tangents are regularly
homotopic. It was used at two places and the peer broke the first: for the
empty support the carrier meets the contact disc in two arcs, not one, and they
cross. Both places now argue instead that the deformation inside the contact
disc is a regular homotopy of a polygonal immersion, because no corner
determinant vanishes along it --- which is how the counterpart's own packs
argue it. The entry is gone from this list.
\end{remark}

\subsection{Citations}\label{sec:axcitations}

The five axiom entries in this subsection are statements taken from the
literature. Their sources are:
\begin{itemize}
\item Axiom~\ref{ax:homfly}: W.~B.~R. Lickorish and K.~C. Millett,
\emph{A polynomial invariant of oriented links}, Topology \textbf{26}
(1987), 107--141, Section~1, especially p.~120, for the complete proof of
existence and uniqueness; and S.~Chmutov and M.~Polyak, \emph{Elementary
combinatorics of the HOMFLYPT polynomial}, Int. Math. Res. Not. IMRN 2010,
no.~3, 480--495, Remarks~1--2 on p.~3 and the proof in Section~5, for the
retained knot-parity clause in the exact normalization of
Axiom~\ref{ax:homfly}.
\item Axiom~\ref{ax:etnyre}: H.~Geiges, \emph{An Introduction to Contact
Topology}, Cambridge Studies in Advanced Mathematics \textbf{109}, Cambridge
University Press (2008), Proposition~3.5.32, p.~127, with printed proof; see
also J.~B. Etnyre, \emph{Legendrian and transversal knots},
\texttt{arXiv:math/0306256}, in \emph{Handbook of Knot Theory}, Elsevier
(2005), 105--185, Section~2.6.4, Equation~(9).
\item Axiom~\ref{ax:slbound}: L. Ng, \emph{A skein approach to Bennequin type
inequalities}, Int. Math. Res. Not. (text at
\texttt{sites.math.duke.edu/\~{}ng/math/papers/bound.pdf}), stated in its list
of Bennequin type inequalities as the HOMFLY--PT bound and proved by its
Theorem~1 (skein template), extended to oriented links by its Corollary~1;
attributed there to J. Franks, R. F. Williams, \emph{Braids and the Jones
polynomial}, Trans. Amer. Math. Soc. \textbf{303} (1987), 97--108, and
H. R. Morton, \emph{Seifert circles and knot polynomials}, Math. Proc.
Cambridge Philos. Soc. \textbf{99} (1986), 107--109; the reading of that bound
at a single transverse representative uses J.~B.~Etnyre, \emph{Legendrian and
transversal knots}, preprint \texttt{arXiv:math/0306256}, Section~2.9 --- the
Legendrian pushoff of a transverse knot, and Lemma~2.22,
$\operatorname{sl}(T_\pm(L))=\operatorname{tb}(L)\mp r(L)$. For the classical
origin of the inequality, rather than the HOMFLY--PT refinement, see
D.~Bennequin, \emph{Entrelacements et \'equations de Pfaff}, Ast\'erisque
\textbf{107--108} (1983), 87--161, Th\'eor\`eme~3, p.~96.
\item Axiom~\ref{ax:bem}: A. Brini, B. Eynard, M. Mari\~no, \emph{Torus knots
and mirror symmetry}, Ann. Henri Poincar\'e \textbf{13} (2012), 1873--1910;
preprint \texttt{arXiv:1105.2012}, the finite closed formula of Section~3 with
the quantum integers of its Equation~(3.40). Read at the preprint on
2026-08-23. The proof source for the coefficient is E.~Gorsky,
\emph{$q,t$-Catalan numbers and knot homology},
\texttt{arXiv:1003.0916v3}, Theorem~3.1 and its Appendix. Its torus-knot
HOMFLY input is V.~F.~R. Jones, \emph{Hecke algebra representations of braid
groups and link polynomials}, Ann. of Math. \textbf{126} (1987), 335--388,
Theorem~9.7, p.~359, whose proof is printed. The rewriting uses N.~M.
Dunfield, S.~Gukov and J.~Rasmussen, \emph{The superpolynomial for knot
homologies}, Experiment. Math. \textbf{15} (2006), 129--159,
Equations~(69) and~(71), p.~25; the equality of those displays, asserted there
without the intervening algebra, has been verified term by term in a symbolic
derivation on file.
\item Axiom~\ref{ax:gausscode}: J.~S.~Carter, \emph{Classifying immersed
curves}, Proc.\ Amer.\ Math.\ Soc.\ \textbf{111} (1991) 281--287, which
classifies generic immersed curves in the plane by their Gauss data up to
homeomorphism of the sphere; the use of Gauss codes as a complete presentation
of a classical link, with the classical codes exactly the realizable ones, is
the framing of L.~H.~Kauffman, \emph{Virtual knot theory}, European J.\
Combin.\ \textbf{20} (1999) 663--690. For diagrams satisfying the paper's
standing Rule~1, C.~H. Dowker and M.~B. Thistlethwaite, \emph{Classification
of knot projections}, Topology Appl. \textbf{16} (1983), 19--31,
Corollary~1.2, p.~24, supplies the orientation refinement; this citation is
only for that prime/reduced subclass. The universal signed-oriented
Gauss-data route uses M.~Goussarov, M.~Polyak and O.~Viro, \emph{Finite type
invariants of classical and virtual knots}, Topology \textbf{39} (2000),
1045--1068, Theorem~1.A, p.~4, together with G.~Kuperberg, \emph{What is a
virtual link?}, Algebr. Geom. Topol. \textbf{3} (2003), 587--591, Theorem~1,
p.~588, and the oriented-link conclusion on p.~591. For the isotopy step see
S.~Smale, \emph{Diffeomorphisms of the 2-sphere}, Proc. Amer. Math. Soc.
\textbf{10} (1959), 621--626, Theorem~6; the topological case is attributed
there to H.~Kneser, \emph{Die Deformationss\"atze der einfach
zusammenh\"angenden Fl\"achen}, Math. Z. \textbf{25} (1926), 362--372.
\end{itemize}
